\section*{第 4 章：算术表达式的 ML 实现}
\addcontentsline{toc}{section}{第 4 章：算术表达式的 ML 实现}

\begin{taplsolutionentrybox}
\phantomsection
\label{sol:translator:4.2.2}
\textbf{练习 4.2.2}

大步求值器直接递归求出各前提的最终值，而不是先生成一个中间项、再从头调用
\texttt{eval}：

\begin{taplocamlcode}
let rec eval t =
  match t with
    TmTrue _ -> t
  | TmFalse _ -> t
  | TmZero _ -> t
  | TmIf(_, t1, t2, t3) ->
      (match eval t1 with
         TmTrue _ -> eval t2
       | TmFalse _ -> eval t3
       | _ -> raise NoRuleApplies)
  | TmSucc(fi, t1) ->
      let v1 = eval t1 in
      if isnumericval v1 then TmSucc(fi, v1)
      else raise NoRuleApplies
  | TmPred(_, t1) ->
      (match eval t1 with
         TmZero _ -> TmZero(dummyinfo)
       | TmSucc(_, nv1) when isnumericval nv1 -> nv1
       | _ -> raise NoRuleApplies)
  | TmIsZero(_, t1) ->
      (match eval t1 with
         TmZero _ -> TmTrue(dummyinfo)
       | TmSucc(_, nv1) when isnumericval nv1 ->
           TmFalse(dummyinfo)
       | _ -> raise NoRuleApplies)
\end{taplocamlcode}
\taplrustcounterpart[break]{sol-translator-04-eval-big}

条件式只求值被守卫选中的分支，这与练习 3.5.17 的大步规则一致。若某个前提
求值到不符合规则所需形状的范式，就抛出 \texttt{NoRuleApplies}，表示原项
卡住。

\taplsolutionbacklink{ex:4.2.2}{返回练习 4.2.2}
\end{taplsolutionentrybox}
